\documentclass[11pt,a4paper]{article}
\usepackage[margin=2.5cm]{geometry}
\usepackage[utf8]{inputenc}
\usepackage[T1]{fontenc}
\usepackage[french]{babel}
\usepackage{hyperref}

\begin{document}
	\pagestyle{empty}
	
	\noindent
	Friedly WOLI \\
	Toulouse, 31100 \\
	0767339290 \\
	friedlysedjro@gmail.com 
	
	\begin{flushright}
		Capgemini Engineering \\
		Toulouse
	\end{flushright}
	
	\bigskip
	
	\textbf{Objet :} Candidature – Ingénieur Simulation Flight Physics \\
	
	Madame, Monsieur,
	
	Actuellement étudiant en mécanique énergétique à l’Université de Toulouse, je souhaite rejoindre votre équipe « Flight Physics » afin de mettre en pratique mes compétences en modélisation et simulation numérique dans le domaine aéronautique.
	
	Au cours de ma formation, j’ai acquis de solides bases en mécanique des fluides, mécanique des solides, transferts thermiques ainsi qu’en méthodes numériques. Ces enseignements m’ont permis de développer une approche rigoureuse de la modélisation physique et de l’analyse des phénomènes complexes. J’ai également développé un fort intérêt pour la simulation numérique appliquée aux problématiques industrielles, notamment en aérodynamique et en performance des systèmes.
	
	Mon expérience au sein de la Toulouse Racing Division à l'ISAE-Supaero illustre ma capacité à appliquer ces concepts à des problématiques concrètes. J’ai notamment participé à la conception d’un véhicule électrique en assurant le codage sous Matlab et Simulink de son fonctionnement global. Ce projet m’a permis de présenter nos résultats lors de la compétition à Silverstone, devant un jury international, en tant que rapporteur de groupe, ce qui a renforcé mes capacités de communication technique en anglais (niveau B2) et mon esprit d'équipe.
	Je maîtrise les outils numériques indispensables tels que Python et Matlab, et possède des compétences sur des logiciels de référence comme Ansys, Abaqus et 3D Experience. Ma curiosité et ma persévérance me permettent de m'adapter rapidement à des méthodes de calcul variées, qu'elles soient analytiques ou basées sur des modèles numériques avancés.
	Intégrer Capgemini Engineering représente pour moi une opportunité de participer à des projets innovants dans le secteur aéronautique et de contribuer au développement des aéronefs de demain, tout en évoluant au sein d’un environnement stimulant et orienté R\&D.
	
	Je me tiens à votre disposition pour tout entretien et vous remercie par avance de l’attention portée à ma candidature.
	
	Je vous prie d’agréer, Madame, Monsieur, l’expression de mes salutations distinguées.
	
	\vspace{0.8cm}
	
	\begin{flushright}
		\textbf{Friedly WOLI}
	\end{flushright}
	
\end{document}